\documentclass[10pt,a4paper]{amsart}
\usepackage[margin=19mm]{geometry}
\usepackage{amsmath,amssymb,mathtools}
\usepackage[scr=rsfs,cal=euler]{mathalfa}
\usepackage{xfrac}
\usepackage[unicode,hidelinks,bookmarks=false]{hyperref}
\newcommand{\ie}{i.e.\ }
\newcommand{\cf}{cf.\ }
\newcommand{\resp}{respectively\ }
\newcommand{\Subsection}{\subsection}
\providecommand{\nobreakdash}{\nobreak\hbox{-}}
\setlength{\emergencystretch}{2em}
\multlinegap0pt
\usepackage{amssymb}
\usepackage{xcolor}
\usepackage{algorithm}\usepackage{algpseudocode}
\usepackage[normalem]{ulem}
\newtheorem{lemma}{Lemma}
\newcommand{\ee}{\mathbb{E}}
\newcommand{\ind}{\mathbb{I}}
\newcommand{\lt}{\left}
\newcommand{\pp}{\mathbb{P}}
\newcommand{\rt}{\right}
\title{Efficient Rare-Event Simulation for Random Geometric Graphs via Importance Sampling}
\author{Sarat Moka, Christian Hirsch, Volker Schmidt,, Dirk Kroese}
\date{Source: arXiv:2504.10530v3 (CC BY 4.0)}
\begin{document}
\maketitle

\noindent\emph{Source and licence.} Efficient Rare-Event Simulation for Random Geometric Graphs via Importance Sampling. Version arXiv:2504.10530v3 (CC BY 4.0), distributed by arXiv under the Creative Commons Attribution 4.0 International licence, http://creativecommons.org/licenses/by/4.0/, as recorded in the arXiv licence metadata for this exact version and shown on the abstract page https://arxiv.org/abs/2504.10530v3. Full licence notice: \texttt{licenses/arxiv-2504.10530-CC-BY-4.0.txt}.\par\smallskip

\noindent\textit{Editorial scope.} Counted unit: the article's lemma lem:upper-bound-ELX (source0.tex lines 1385--1397). The article's complete original proof of it is delivered with it (source0.tex lines 1398--1460, one \texttt{proof} environment of 63 lines). No mathematics is written, repaired or re-derived: the statement and the proof are the article's own lines, in the article's own notation, and no retained line is substituted or re-flowed.  No further source-local context is needed: every cross-reference of every retained line resolves inside the two delivered spans.  Nothing was omitted from any retained span, so the package carries no omission record.  No retained line contains a citation, so the wrapper carries no bibliography and no bibliography entry.  No retained line contains a figure environment, an image-inclusion command or drawing code, so no image is delivered and the package declares no asset dependency.  Editorial formatting only, in addition: an AMS \texttt{amsart} preamble that restates the article's own declarations and macros used by the retained lines, namely the environments \texttt{lemma} on separate counters, together with the standard packages those lines need.  The retained environments print in this document as Lemma 1 in order of appearance, whereas the article prints the same items with the numbering of its own full document.  The complete native source of the article as distributed in the arXiv e-print for this exact version is bundled unchanged as \texttt{sources/176532/source0.tex}.
\begin{lemma}
    \label{lem:upper-bound-ELX}
    For any $n \leq 2\beta$, 
    \begin{align}
    \sqrt{\ee_\mu\lt[\exp\lt(\frac{2}{\beta^{\delta}}  \sum_{i =1}^{n-1} (i\upsilon_d - V_i)\rt)\rt]} = \exp\lt( O\lt(\beta^{2 -{\delta(d+1)/d}} + \beta^{3 - 2\delta} \rt)\rt),
    \label{eqn:part1-of-lemma}
\end{align}
and 
\begin{align}
 \sum_{i = 1}^{n -1} \sum_{j = 2}^\infty\frac{(i\upsilon_d)^j}{\beta^{j\delta}} = O( \beta^{3 - 2\delta}). 
\label{eqn:part2-of-lemma}
\end{align}
\end{lemma}

\begin{proof}
For each $i = 1, \dots, n-1$, define the indicator random variable
$I_i = \ind( V_{i} - V_{i-1} < \upsilon_d)$.
Then, 
\begin{align}
i\upsilon_d - V_{i} &= \sum_{j = 1}^{i} \lt(\upsilon_d - (V_{j} - V_{j-1}) \rt) =  \sum_{j =1}^i I_i\lt(\upsilon_d - (V_{j} - V_{j-1}) \rt) 
\leq {  \upsilon_d \sum_{j =1}^i I_j,}
\label{eqn:low_bdd_Bi}
\end{align}
Thus, 
\begin{align*}
\sum_{i =1}^{n-1} (i\upsilon_d - V_i) \leq \upsilon_d \sum_{i = 1}^{n-1} \sum_{j = 1}^i I_j = \upsilon_d \sum_{j = 1}^{n-1} (n - j) I_j,
\end{align*}
for each $i = 1,\dots, n-1$.
Note that $I_i = 1$ means the $i^{th}$ point $X_i$ is generated within $1$ unit from the boundary of the window~$W$ or within $2$ units from the centers of all the existing $i-1$ points. {  Observe that  for any $j =1, \dots, i-1$, the probability of $\|X_i - X_j\|_2 \leq 2$ is at most $2^d \upsilon_d/(\beta^{\delta} - V_i)$. This is because, $X_j$ is generated on non-blocking region created by the points $1, \dots, i-1$ and $x_i$ has to be within two units distance from $X_i$.} Thus, we have
\[
\pp_\mu \lt( I_i = 1 | X_1,\dots, X_{i-1}\rt) \leq 
\frac{\beta^{\delta} - (\beta^{\delta/d} - 2)^d}{\beta^{\delta} - V_i} + \sum_{j = 1}^{i-1} \frac{2^d \upsilon_d}{\beta^{\delta} - V_i}.
\]
Since $V_i \leq i \upsilon_d,$ with $r_i =  \frac{\beta^{\delta} - (\beta^{{\delta/d} } -2)^d}{\beta^{\delta} - i\upsilon_d} + \frac{i \upsilon_d}{\beta^{\delta} - i \upsilon_d}$, we have 
$\pp_\mu\lt( I_i = 1| X_1,\dots, X_{i-1}\rt) \leq r_i.$
Consequently,
\begin{align*}
\ee_\mu\lt[\exp\lt(\frac{2}{\beta^{\delta}}  \sum_{i =1}^{n-1} (i\upsilon_d - V_i)\rt)\rt] &\leq \ee_\mu\lt[\exp\lt(\frac{2 \upsilon_d}{\beta^{\delta}} \sum_{j = 1}^{n-1} (n - j) I_j \rt)\rt].
\end{align*}
Observe that the right side of the above inequality is equal to 
\begin{align*}
{  \ee_\mu\lt[\exp\lt(\frac{2 \upsilon_d}{\beta^{\delta}} \sum_{j = 1}^{n-2} (n - j) I_j \rt) \ee_\mu\lt[ \exp\lt(\frac{2 \upsilon_d}{\beta^{\delta}} I_{n-1} \rt)\big| X_1,\dots, X_{n-2}\rt]\rt]},
\end{align*}
which can be upper bounded by
\begin{align*}
\ee_\mu\lt[\exp\lt(\frac{2 \upsilon_d}{\beta^{\delta}} \sum_{j = 1}^{n-2} (n - j) I_j \rt) \rt] \lt((1 - r_{n-1}) + r_{n-1} \exp(2 \upsilon_d/\beta^{\delta}) \rt).
\end{align*}
By repeating the same procedure for all the terms in the above expectation and substituting the values of $c_i$'s, we establish that 
\begin{align}
\ee_\mu\lt[\exp\lt(\frac{2}{\beta^{\delta}}  \sum_{i =1}^{n-1} (i\upsilon_d - V_i)\rt)\rt] 
&\leq \prod_{i = 1}^{n-1}\lt(1 + r_{n-i} (e^{2i \upsilon_d/\beta^{\delta}} -1) \rt)\nonumber\\
&\leq \exp\lt( \sum_{i = 1}^{n-1} r_{n-i} (e^{2i \upsilon_d/\beta^{\delta}} -1) \rt). 
\label{eqn:exp-exp-bdd}
\end{align}
Then, notice that since $\delta > 1$, for each $i < n \leq 2\beta$, $iv/\beta^{\delta} \to 0$ as $\beta \to \infty$. Furthermore, 
 $\beta^{\delta} - (\beta^{\delta/d} - 2)^d = O(\beta^{\delta(d -1)/d})$, and hence,  
$r_i = O\lt((\beta^{\delta(d-1)/d} + i)/\beta^{\delta}\rt)$.
Also,
\begin{align*}
    \exp(2i \upsilon_d/\beta^{\delta}) -1 = \sum_{j=1}^\infty \frac{(2i \upsilon_d)^j}{j!\beta^{j\delta}} = \frac{2i\upsilon_d}{\beta^{\delta}} \sum_{j=0}^\infty \frac{(2i \upsilon_d)^j}{(j+1)!\beta^{j\delta}} \leq \frac{2i\upsilon_d}{\beta^{\delta}} \exp(2i \upsilon_d/\beta^{\delta}).
\end{align*}
Since $\exp(2i \upsilon_d/\beta^{\delta}) \to 1$ as $\beta \to \infty$ for all $i < n \leq 2\beta$, we have
\[
\sum_{i = 1}^{n-1} r_{n-i} (\exp(2i \upsilon_d/\beta^{\delta}) -1) = O\lt(\frac{n^2\beta^{\delta(d-1)/d}}{\beta^{2\delta}} + \frac{n^3}{\beta^{2\delta}} \rt) = O\lt(\beta^{2 - {\delta(d+1)/d}} + \beta^{3 - 2\delta} \rt),
\]
which completes the proof of \eqref{eqn:part1-of-lemma} using \eqref{eqn:exp-exp-bdd}.
Now to prove \eqref{eqn:part1-of-lemma}, observe for large $\beta$ that if $i < n \leq 2\beta$, 
\begin{align*}
    \sum_{i = 1}^{n -1} \sum_{j = 2}^\infty \frac{(i\upsilon_d)^j}{\beta^{j\delta}} 
    = \frac{\upsilon_d^2}{\beta^{2\delta}} \sum_{i = 1}^{n -1} i^2 \sum_{j = 0}^\infty \frac{(i\upsilon_d)^j}{\beta^{j\delta}} 
    = \frac{\upsilon_d^2}{\beta^{2\delta}} \sum_{i = 1}^{n -1} i^2 \frac{1}{1 - i\upsilon_d/\beta^{\delta}}
    \leq \frac{\upsilon_d^2}{\beta^{2\delta}} \frac{1}{1 - n\upsilon_d/\beta^{\delta}}\sum_{i = 1}^{n -1} i^2,
\end{align*}
where we used $i\upsilon_d < n\upsilon_d$ in the last inequality.
We complete the proof by observing that $n/\beta^\delta \leq 2\beta/\beta^\delta \to 0$ as $\beta \to \infty$ and 
$\sum_{i < n} i^2 \leq \sum_{i < 2\beta} i^2 = O(\beta^3).$
\end{proof}

\end{document}
